\documentclass[10pt,a4paper]{amsart}
\usepackage{amsmath,amssymb,amsthm,graphicx}
% Retained minimal preamble for the Stacks Project chapter more-algebra.tex (More on Algebra).
% The chapter file pinned-original-source.tex carries only \input{preamble} on its first line, so no
% preamble text lives inside the chapter itself.  The Stacks Project uses one single shared preamble for
% every chapter; the snapshot stages it identically as sources/stacks-chow/preamble.tex and
% sources/stacks-etale-cohomology/preamble.tex.  The block below is copied mechanically, line for line,
% from sources/stacks-chow/preamble.tex lines 14-18 (the Xy-pic setup) and 154-217 (the native theorem
% declarations together with the mathematical macros).  Nothing was added, removed or redefined; the
% wrapper class and the amsmath/amssymb/amsthm/graphicx base packages are supplied by the assembler.
% Chapter: more-algebra.tex -- More on Algebra.
% For commutative diagrams we use Xy-pic
\usepackage[all]{xy}

% We use 2cell for 2-commutative diagrams.
\xyoption{2cell}

% Theorem environments.
%
\theoremstyle{plain}
\newtheorem{theorem}[subsection]{Theorem}
\newtheorem{proposition}[subsection]{Proposition}
\newtheorem{lemma}[subsection]{Lemma}

\theoremstyle{definition}
\newtheorem{definition}[subsection]{Definition}
\newtheorem{example}[subsection]{Example}
\newtheorem{exercise}[subsection]{Exercise}
\newtheorem{situation}[subsection]{Situation}

\theoremstyle{remark}
\newtheorem{remark}[subsection]{Remark}
\newtheorem{remarks}[subsection]{Remarks}

\numberwithin{equation}{subsection}

% Macros
%
\def\lim{\mathop{\mathrm{lim}}\nolimits}
\def\colim{\mathop{\mathrm{colim}}\nolimits}
\def\Spec{\mathop{\mathrm{Spec}}}
\def\Hom{\mathop{\mathrm{Hom}}\nolimits}
\def\Ext{\mathop{\mathrm{Ext}}\nolimits}
\def\SheafHom{\mathop{\mathcal{H}\!\mathit{om}}\nolimits}
\def\SheafExt{\mathop{\mathcal{E}\!\mathit{xt}}\nolimits}
\def\Sch{\mathit{Sch}}
\def\Mor{\mathop{\mathrm{Mor}}\nolimits}
\def\Ob{\mathop{\mathrm{Ob}}\nolimits}
\def\Sh{\mathop{\mathit{Sh}}\nolimits}
\def\NL{\mathop{N\!L}\nolimits}
\def\CH{\mathop{\mathrm{CH}}\nolimits}
\def\proetale{{pro\text{-}\acute{e}tale}}
\def\etale{{\acute{e}tale}}
\def\QCoh{\mathit{QCoh}}
\def\Ker{\mathop{\mathrm{Ker}}}
\def\Im{\mathop{\mathrm{Im}}}
\def\Coker{\mathop{\mathrm{Coker}}}
\def\Coim{\mathop{\mathrm{Coim}}}

% Boxtimes
%
\DeclareMathSymbol{\boxtimes}{\mathbin}{AMSa}{"02}

%
% Macros for moduli stacks/spaces
%
\def\QCohstack{\mathcal{QC}\!\mathit{oh}}
\def\Cohstack{\mathcal{C}\!\mathit{oh}}
\def\Spacesstack{\mathcal{S}\!\mathit{paces}}
\def\Quotfunctor{\mathrm{Quot}}
\def\Hilbfunctor{\mathrm{Hilb}}
\def\Curvesstack{\mathcal{C}\!\mathit{urves}}
\def\Polarizedstack{\mathcal{P}\!\mathit{olarized}}
\def\Complexesstack{\mathcal{C}\!\mathit{omplexes}}
% \Pic is the operator that assigns to X its picard group, usage \Pic(X)
% \Picardstack_{X/B} denotes the Picard stack of X over B
% \Picardfunctor_{X/B} denotes the Picard functor of X over B
\def\Pic{\mathop{\mathrm{Pic}}\nolimits}
\def\Picardstack{\mathcal{P}\!\mathit{ic}}
\def\Picardfunctor{\mathrm{Pic}}
\def\Deformationcategory{\mathcal{D}\!\mathit{ef}}

% --- appended 2026-10-12 (Chief): the project own definitions of the reference/slogan/history
% display environments, copied mechanically from the shared Stacks preamble lines 8-12
% (\usepackage{verbatim} plus the three \newenvironment lines, which the chapter preamble keeps
% verbatim from the project preamble). Nothing added, removed or redefined: these environments
% carry no mathematical content and the project itself typesets them as comments. Without these
% three lines a counted statement whose own text uses them fails with Environment slogan undefined.
\usepackage{verbatim}
\newenvironment{reference}{\comment}{\endcomment}
\newenvironment{slogan}{\comment}{\endcomment}
\newenvironment{history}{\comment}{\endcomment}
\usepackage{cleveref}
\crefname{originallemma}{original lemma}{original lemmas}
\Crefname{originallemma}{Original lemma}{Original lemmas}
\makeatletter
\newlabel{lemma-derived-completion}{{091V}{source}{}{sourcepointer-1}{}}
\newlabel{lemma-derived-completion@cref}{{[originallemma][0][]091V}{[1][0][]source}{}{}{}}
\makeatother
\crefname{originallemma}{original lemma}{original lemmas}
\Crefname{originallemma}{Original lemma}{Original lemmas}
\makeatletter
\newlabel{lemma-derived-completion-finite-cohomological-dimension}{{0AAJ}{source}{}{sourcepointer-2}{}}
\newlabel{lemma-derived-completion-finite-cohomological-dimension@cref}{{[originallemma][0][]0AAJ}{[1][0][]source}{}{}{}}
\makeatother
\crefname{originallemma}{original lemma}{original lemmas}
\Crefname{originallemma}{Original lemma}{Original lemmas}
\makeatletter
\newlabel{lemma-projective-amplitude}{{0A5P}{source}{}{sourcepointer-3}{}}
\newlabel{lemma-projective-amplitude@cref}{{[originallemma][0][]0A5P}{[1][0][]source}{}{}{}}
\makeatother
\crefname{originalsection}{original section}{original sections}
\Crefname{originalsection}{Original section}{Original sections}
\makeatletter
\newlabel{homology-section-filtered-complex}{{012K}{source}{}{sourcepointer-4}{}}
\newlabel{homology-section-filtered-complex@cref}{{[originalsection][0][]012K}{[1][0][]source}{}{}{}}
\makeatother
\crefname{originallemma}{original lemma}{original lemmas}
\Crefname{originallemma}{Original lemma}{Original lemmas}
\makeatletter
\newlabel{homology-lemma-biregular-ss-converges}{{012W}{source}{}{sourcepointer-5}{}}
\newlabel{homology-lemma-biregular-ss-converges@cref}{{[originallemma][0][]012W}{[1][0][]source}{}{}{}}
\makeatother
\title{A canonical spectral sequence of derived complete modules computing the completion of a filtered complex}
\author{The Stacks Project Authors}
\date{stacks-project:more-algebra.tex:local-snapshot-20261012; GNU FDL 1.2 or later}
\begin{document}
\maketitle
\setcounter{subsection}{1}
\noindent\textit{Editorial note.} All delivered mathematical text is the author's own, copied from the pinned Stacks Project chapter file more-algebra.tex; no statement or proof step was rewritten, reordered or removed. The counted statement is the labeled lemma environment at source lines 26555-26567 (canonical tag 0BKD) and the counted proof is the proof environment at source lines 26569-26591, which immediately follows it and is its own proof.\par\medskip
\noindent\textit{Editorial note.} The statement and proof use the notation of graded pieces and of the completion of a complex from the surrounding section; the completion functor is the one whose description is carried as a navigation stub.\par\medskip
\noindent\textit{Editorial note.} The preamble span runs to line 90, the current end of the retained preamble, which includes the appended environment definitions; the counted text does not use them, and the span is taken whole so that the file version is unambiguous.\par\medskip
\noindent\textit{Editorial note.} The source URL is the source-level Stacks Project URL, as required when more than one card is produced from one source; the per-claim tag is carried by claim\_key and by this note. The counted statement and proof cite no bibliography entry, so this card needs no bibliography.\par\medskip
\noindent\textit{Allowed standard prerequisites (editorial).}\begin{itemize}
\item Derived completion with respect to a finitely generated ideal is given by the derived Hom from a fixed bounded complex, and by the finite cohomological dimension statement that complex may be taken of finite projective dimension. invoked at proof lines 26569-26573. Replaces the completion functor by a Hom complex from a bounded complex of projectives, which is the object whose filtration is studied.
\item The Hom complex from a bounded complex of projectives into a complex represents the derived Hom, and applying it to a filtered complex produces a filtered complex with the induced filtration. used at proof lines 26574-26581. Produces the complex representing the completion together with the filtration whose spectral sequence is the one claimed.
\item The spectral sequence of a filtered complex has the stated first page and differential bidegree, and for a finite filtration on each term it is bounded and converges to the cohomology of the whole complex. invoked at proof lines 26581-26590. Gives the identification of the first page with the cohomology of the completion of the graded pieces and the convergence statement.
\end{itemize}\medskip
\ifdefined\hypertarget\hypertarget{sourcepointer-1}{}\fi
\noindent\textit{Source reference.} Original lemma 091V: sources/stacks-more-algebra/pinned-original-source.tex, source lines 26154--26154. The reference is invoked at proof line 26569 to write derived completion as a derived Hom from a fixed object; the construction of the Hom complex and its filtration are inside the counted span.\par\medskip
\ifdefined\hypertarget\hypertarget{sourcepointer-2}{}\fi
\noindent\textit{Source reference.} Original lemma 0AAJ: sources/stacks-more-algebra/pinned-original-source.tex, source lines 26507--26507. The reference together with the amplitude statement gives finite projective dimension of the representing object, used at proof lines 26570-26573; everything built on it is inside the counted span. The lemma itself is the subject of a separate card of this lane with disjoint coordinates.\par\medskip
\ifdefined\hypertarget\hypertarget{sourcepointer-3}{}\fi
\noindent\textit{Source reference.} Original lemma 0A5P: sources/stacks-more-algebra/pinned-original-source.tex, source lines 12817--12817. The reference is the second input to the conclusion that the representing object has finite projective dimension, cited at proof lines 26570-26572; it is a standard pointer and the use made of it is inside the counted span.\par\medskip
\ifdefined\hypertarget\hypertarget{sourcepointer-4}{}\fi
\noindent\textit{Source reference.} Original section 012K: sources/stacks-chow/tags.txt, source lines 1402--1402. The reference points to the framework of spectral sequences of filtered complexes, cited at proof line 26582 to obtain the spectral sequence with the stated first page; it is a pointer to the standard construction and the filtration is built inside the counted span.\par\medskip
\ifdefined\hypertarget\hypertarget{sourcepointer-5}{}\fi
\noindent\textit{Source reference.} Original lemma 012W: sources/stacks-chow/tags.txt, source lines 1414--1414. The reference is the convergence criterion applied at proof lines 26588-26590 after the filtration is seen to be finite on each term; the finiteness argument itself is inside the counted span.\par\medskip
\begin{lemma}
\label{lemma-derived-completion-spectral-sequence}
Let $A$ be a ring and let $I \subset A$ be a finitely generated ideal.
Let $K^\bullet$ be a filtered complex of $A$-modules. There exists a
canonical spectral sequence $(E_r, \text{d}_r)_{r \geq 1}$
of bigraded derived complete $A$-modules with $d_r$ of bidegree
$(r, -r + 1)$ and with
$$
E_1^{p, q} = H^{p + q}((\text{gr}^pK^\bullet)^\wedge)
$$
If the filtration on each $K^n$ is finite, then the spectral sequence is
bounded and converges to $H^*((K^\bullet)^\wedge)$.
\end{lemma}

\begin{proof}
By Lemma \ref{lemma-derived-completion} we know that derived completion is
given by $R\Hom_A(C, -)$ for some $C \in D^b(A)$. By
Lemmas \ref{lemma-derived-completion-finite-cohomological-dimension} and
\ref{lemma-projective-amplitude} we see that $C$ has finite projective
dimension. Thus we may choose a bounded complex of projective
modules $P^\bullet$ representing $C$. Then
$$
M^\bullet = \Hom^\bullet(P^\bullet, K^\bullet)
$$
is a complex of $A$-modules representing $(K^\bullet)^\wedge$.
It comes with a filtration given by
$F^pM^\bullet = \Hom^\bullet(P^\bullet, F^pK^\bullet)$.
We see that $F^pM^\bullet$ represents $(F^pK^\bullet)^\wedge$
and hence $\text{gr}^pM^\bullet$ represents $(\text{gr}K^\bullet)^\wedge$.
Thus we find our spectral sequence by taking the spectral sequence of
the filtered complex $M^\bullet$, see
Homology, Section \ref{homology-section-filtered-complex}.
If the filtration on each $K^n$ is finite, then the filtration
on each $M^n$ is finite because $P^\bullet$ is a bounded complex.
Hence the final statement follows from
Homology, Lemma \ref{homology-lemma-biregular-ss-converges}.
\end{proof}

\end{document}
